\chapter{Data and panel construction}
\label{ch:data}

\section{The FNSPID panel}

FNSPID combines about 15.7 million public financial-news records with 29.7
million US equity-price observations \citep{dong2024}. I retain firms with
coverage in every calendar year.

The fixed rule requires at least 20 mappable news days in every calendar year
from 2011 through 2023. It returns 574 ticker strings. Three lack price files and
one produces no joined firm-day, leaving 570 after joining news, adjusted opens
and SPY prices. SPY is an exchange-traded fund that tracks the S\&P 500 and is
used here as the market benchmark. This coverage-filtered firm universe favours
surviving, well-covered firms. It follows data availability rather than
historical index membership and applies no further security-type filter.

The panel contains 715,546 news-bearing firm-days across 3,262 sessions:
512,153 in training and 203,393 in testing. Four training outcomes are missing,
so 512,149 rows enter the nine-rule comparison.

Training-period coverage is sparse: 247,470 firm-days (48.3\%) have one story,
127,541 (24.9\%) have two, and 137,142 (26.8\%) have at least three. On a
one-story day, negative share is a binary label rather than a distribution
across stories.

Each story row stores its symbol, assigned session, whether it is marked as a
recap and three FinBERT class probabilities. Rows are collapsed into the
firm-session counts, means,
dispersion, label shares and other rules in \cref{ch:methodology}. Each
observation combines one priced firm with one assigned exchange session; this is
the analytic unit $(i,t)$. Story text remains in a local database and never
enters this repository or its intermediate files.

Timing follows the fixed session-mapping rule. Date-only records move to the first
New York Stock Exchange session after their date, whereas timestamped records
keep the session implied by that time. The date-only rule can therefore delay pre-open
news by a full session. Because the data contain neither verified availability
times nor a row-level timing-precision flag, it is not possible to identify the
affected rows.
Some stories may have appeared after their assigned return window opened, which
prevents an interpretation based on news known before the return began.

\section{Scoring a story}

FinBERT assigns each story a probability of positive, neutral and negative
sentiment. The same fixed ProsusAI FinBERT classifier scores both data sources.
They use model identifier \texttt{ProsusAI/finbert} at revision
\texttt{4556d13015211d73\allowbreak dccd3fdd39d39232506f3e43}, loaded from an
enforced local snapshot. For story $j$ about firm $i$ on assigned session $t$,
the continuous score is

\begin{equation}
 s_{ijt}=p_{ijt}(\mathrm{positive})-p_{ijt}(\mathrm{negative}),
 \qquad -1\leq s_{ijt}\leq 1.
 \label{eq:continuous-score}
\end{equation}

The hard label is the class with the highest probability. A negative probability
of 0.40 wins if the other two are lower, so the boundary is relative rather than
a confidence threshold. The model was pretrained on financial text and
fine-tuned on Financial PhraseBank \citep{araci2019}. FNSPID has no human-label
audit or independent-model comparison, so the analysis concerns the unverified
outputs of one fixed classifier.

The mix of sentiment classes changes modestly at the period boundary. The
1,183,961 training
stories are 16.80\% negative, 60.01\% neutral and 23.20\% positive. The 415,490
testing stories are 15.84\%, 61.42\% and 22.74\%. Negative share is zero on
74.93\% of training firm-days and 76.85\% of testing firm-days. Single-story
days rise from 48.32\% to 56.53\%. These shifts do not explain the sign reversal
alone. Timing, mapping, coverage, classifier behaviour and composition may also
matter.

Training covers 2,264 sessions from 3 January 2011 to 31 December 2019. Testing
covers 998 sessions from 2 January 2020 to 18 December 2023. Before the testing
estimates were calculated, the two coefficients, multiple-testing adjustment,
power assumptions and decision rule were fixed. The test then ran once without
changes.

Some later portfolio outcomes had already been examined. This decision was
recorded internally rather than in an independently timestamped
preregistration, so the testing period was not completely unseen.

\section{The return being explained}

Daily prices come from the FNSPID archive, whose reference collection code
downloads Yahoo Finance bars. Yahoo publishes an adjusted close but not an
adjusted open, so I build the proxy\footnote{See the \href{https://github.com/Zdong104/FNSPID_Financial_News_Dataset/blob/main/data_scraper/stock_price_scraper/get_price_from_yahoo.py}{FNSPID price scraper} and \href{https://help.yahoo.com/kb/adjusted-close-sln28256.html}{Yahoo adjusted-close documentation}, both accessed 15 August 2026.}

\begin{equation}
 P^{\mathrm{adj,open}}_{it}=P^{\mathrm{open}}_{it}
 \frac{P^{\mathrm{adj,close}}_{it}}{P^{\mathrm{close}}_{it}}.
 \label{eq:adjusted-open}
\end{equation}

Applying Yahoo's adjusted-close ratio to the open transfers its split and
dividend adjustments. I did not reconstruct corporate actions independently.
The result is therefore neither an audited total-return series nor an exact
implementation price. Prices advance to the next SPY-calendar session, not the
firm's next news day, avoiding multi-session returns when an intervening session
has no story. The primary outcome is

\begin{equation}
 r^{A}_{i,t+1}=\left(\frac{P^{\mathrm{adj,open}}_{i,t+1}}
 {P^{\mathrm{adj,open}}_{it}}-1\right)-
 \left(\frac{P^{\mathrm{adj,open}}_{\mathrm{SPY},t+1}}
 {P^{\mathrm{adj,open}}_{\mathrm{SPY},t}}-1\right).
 \label{eq:abnormal-return}
\end{equation}

Subtracting SPY's open-to-open return gives a simple market-adjusted event-study
outcome \citep{mackinlay1997}. Two additional sensitivity checks replace unit market
exposure with a 252-session trailing stock-SPY beta. One beta-adjusts the
outcome; the other also controls for beta rank. The baseline association and
period contrast persist, but remain rank associations, not estimates of factor
alpha.

\section{The Reuters/LSEG data}

Reuters coverage supplied through LSEG covers 33 large US firms across several
sectors. It is denser and has richer source and availability fields than the
available FNSPID data, making it useful as a separate second sample.

The two non-overlapping LSEG periods both follow the FNSPID training period. The
earlier period has 14,935 firm-days across 456 sessions (2 January 2024--24
October 2025). The later period has 5,511 rows across 167 sessions (27 October
2025--26 June 2026). Both cover 33 firms.

News density prevents a direct replication. Median daily story count is 61 in
the earlier period and 90 in the later period, versus 2 in FNSPID training. A
large share change may flip one of two FNSPID labels, but shifts a dense
distribution in LSEG.

LSEG maps each story to the first New York Stock Exchange (XNYS) open at least
fifteen minutes after recorded availability, which FNSPID cannot do. The
cross-source table uses raw open-to-open returns because the LSEG aggregate lacks
the FNSPID abnormal-return construction. A common centred percentile-rank
transform puts the coefficients on a comparable scale. The outcomes remain
different and are labelled accordingly. The sources are never pooled.

The firm-level LSEG economic test requires all 33 firms to be present in every
included session. This removes 11 sessions from the earlier period. The complete
panel contains 612 sessions: 445 for training and 167 for testing. The training
sessions are divided into three chronological folds without changing their time
order. A single rule selected from these folds is then tested on the later
period.

The rule is compared with an equal-weight long portfolio and two
constant-exposure controls selected in training.

Licensed source records, complete request payloads and raw model responses are
excluded from this repository. The prompt appendix reproduces the available
templates and describes the headline fields sent to the model provider.

\section{Data limitations}

FNSPID lacks stable publisher identifiers, near-duplicate clusters and links
from an update to its first release. Several negative rows may therefore restate
one event. Controlling for
$\log(1+n_{it})$ separates score patterns from coverage volume, but not new
information from repetition, a distinction that affects market response
\citep{tetlock2011stale}.

Source density differs sharply, classifier errors remain possible, and FNSPID
requires news, a mapped symbol, a stock open and a SPY open. Average firms per
session fall from 226.2 in training to 203.8 in testing. There is no
common-universe re-estimate, publisher breakdown or row-level timestamp audit.
The period difference shows that the estimated relationship is not stable across
these periods, but it does not explain why. Market conditions, coverage,
production, mapping and classifier behaviour all remain possible explanations.
